\chapter{A Concrete Naming Certificate for $\CompactUniformModulusUp$}
\label{ch:concrete-instances-compact-uniform-modulus-namecert}
\origin{human}

Following Bishop--Bridges compact metric continuity, the $\CompactUniformModulusUp$ packet realizes the finite-net fold named in \autoref{ch:visions-compact-uniform-continuity}. It consumes the metric source of \autoref{ch:concrete-instances-metric-namecert}, the compact finite-cover surface of \autoref{ch:concrete-instances-compact-namecert}, the pointwise continuity rows of \autoref{ch:concrete-instances-continuous-namecert}, the output interface of \autoref{ch:concrete-instances-uniformmodulus-namecert}, and the rational comparison rows of \autoref{ch:concrete-instances-rat-namecert}. The packet names a finite uniform-modulus ledger, not an open-cover compactness theorem, selected Lebesgue number, sequential compactness principle, or host modulus function.

\begin{definition}[Compact uniform modulus carrier]
\label{def:compact-uniform-modulus-namecert-carrier}
A $\CompactUniformModulusUp$ carrier is a finite $\BHist$ packet
$$
M=(X,K,F,L,Q,U,H,C,P,N)
$$
with the following displayed rows:
\begin{enumerate}
\item $X$ is the source $\MetricUp$ row, including the distance and comparison reads used by the compact and continuity surfaces.
\item $K$ is the $\CompactUp$ finite-net row, displaying the finite center bundle and covering radius ledger over $X$.
\item $F$ is the $\ContinuousUp$ pointwise-modulus row, attaching a source radius to each displayed center and requested target tolerance.
\item $L$ is the local finite-net fold ledger, matching each center from $K$ with the corresponding pointwise row from $F$.
\item $Q$ is the $\RatUp$ comparison row computing the finite lower bound used as the source tolerance.
\item $U$ is the $\UniformModulusUp$ output row produced from the fold ledger and rational comparison.
\item $H$ records componentwise $\hsame$ transport for the metric, compact, continuity, fold, rational, output, replay, provenance, and local naming rows.
\item $C$ records displayed $\Cont$ replay from a target-tolerance request through $X,K,F,L,Q,U,H,P,N$.
\item $P$ is the $\Pkg$ provenance over $\CompactUniformModulusUp$, $\MetricUp$, $\CompactUp$, $\ContinuousUp$, $\UniformModulusUp$, $\RatUp$, $\BHist$, $\hsame$, $\Cont$, and $\NameCert$.
\item $N$ is the local $\NameCert_{\CompactUniformModulusUp}$ row.
\end{enumerate}
The classifier compares accepted carriers only through $X,K,F,L,Q,U,H,C,P,N$. It has no coordinate for selected subcovers, selected Lebesgue numbers, sequence extraction, countable choice, host minima, or ambient completeness.
\end{definition}

\begin{theorem}[Compact uniform modulus NameCert obligations]
\label{thm:compact-uniform-modulus-namecert-obligations}
For an accepted $\CompactUniformModulusUp$ carrier $M=(X,K,F,L,Q,U,H,C,P,N)$, the local $\NameCert_{\CompactUniformModulusUp}$ surface consists exactly of carrier admission for the metric source, compact finite-net row, pointwise continuity row, local fold ledger, rational comparison row, uniform-modulus output, transport, replay, provenance, and local naming; source lock between $X,K,F$; finite coverage of every center named by $K$; rational lower-bound exactness through $Q$; output admission through $U$; and ledger non-escape from selected subcovers, selected Lebesgue numbers, sequence extraction, choice, host minima, and ambient completeness.
\end{theorem}
\begin{proof}
Project the carrier of \autoref{def:compact-uniform-modulus-namecert-carrier}. The source metric row $X$ is shared by the compact row $K$ and the continuity row $F$, so the fold cannot change the underlying metric carrier.

The finite-net fold is displayed by $L$. It pairs each center recorded in $K$ with a pointwise continuity radius from $F$, giving a finite ledger rather than a host-selected subcover or sequence argument.

The rational comparison row $Q$ computes the lower bound consumed by the output row $U$. Transport, replay, provenance, and local naming preserve this route, so the absent compactness and choice principles have no coordinate in the NameCert surface.
\end{proof}

\begin{theorem}[Compact finite net to uniform modulus]
\label{thm:compact-uniform-modulus-finite-net-fold}
Let $M=(X,K,F,L,Q,U,H,C,P,N)$ be an accepted $\CompactUniformModulusUp$ carrier. Every uniform-modulus read exported by $M$ factors through the finite route
$$
X,K,F,L,Q,U,H,C,P,N .
$$
The route first reads the compact finite-net rows and pointwise continuity rows, then folds their displayed rational bounds into $Q$, then exposes the $\UniformModulusUp$ row $U$. It never passes through selected subcovers, selected Lebesgue numbers, sequential compactness, host modulus functions, countable choice, or ambient complete metric spaces.
\end{theorem}
\begin{proof}
The carrier definition displays the only modulus-producing coordinates. A read cannot reach $U$ before it has fixed the common source metric $X$, the finite-net row $K$, and the pointwise continuity row $F$.

The fold ledger $L$ and rational comparison row $Q$ provide the finite lower-bound computation. Thus the output is obtained from named rational data, not from an external minimum or selected Lebesgue number.

Finally $H,C,P,N$ transport, replay, source, and name the same route. Since the excluded principles do not occur among the carrier rows, no accepted uniform-modulus read can factor through them.
\end{proof}

\begin{theorem}[Compact uniform modulus RealSeparability handoff]
\label{thm:compact-uniform-modulus-real-separability-handoff}
A $\RealSeparabilityUp$ or compact-metric consumer may read an accepted $\CompactUniformModulusUp$ carrier $M=(X,K,F,L,Q,U,H,C,P,N)$ only through the finite route
$$
\MetricUp,\ \CompactUp,\ \ContinuousUp,\ \UniformModulusUp,\ \RatUp,\ \BHist,\ \hsame,\ \Cont,\ \Pkg,\ \NameCert .
$$
The read exposes the metric source, compact finite-net row, pointwise continuity row, finite fold ledger, rational comparison row, uniform-modulus output, transport, replay, provenance, and local naming. It does not expose a selected Lebesgue number, selected subcover, host modulus function, sequence extraction principle, or ambient compact-metric completeness theorem.
\end{theorem}
\begin{proof}
Project the carrier of \autoref{def:compact-uniform-modulus-namecert-carrier}. The rows $X,K,F,L,Q,U,H,C,P,N$ give the metric source, compact finite-net row, pointwise continuity row, local fold, rational comparison, output, transport, replay, provenance, and local naming.

Use the obligation surface of \autoref{thm:compact-uniform-modulus-namecert-obligations}. It fixes source lock, finite coverage, rational lower-bound exactness, output admission, and non-escape from selected compactness data as local NameCert obligations.

Finally apply the finite-net fold theorem \autoref{thm:compact-uniform-modulus-finite-net-fold}. The fold forces every exported uniform-modulus read through $X,K,F,L,Q,U,H,C,P,N$, so a RealSeparability or compact-metric consumer receives only the displayed finite route and no selected Lebesgue number or host modulus.
\end{proof}

\begin{closurestatus}{\CompactUniformModulusUp}
  \theoryclosure{\seedClosure}
  \scopeclosed{\autoref{def:compact-uniform-modulus-namecert-carrier}, \autoref{thm:compact-uniform-modulus-namecert-obligations}, and \autoref{thm:compact-uniform-modulus-finite-net-fold} seed the finite uniform-modulus packet over metric, compact finite-net, pointwise continuity, local fold, rational comparison, output, transport, replay, provenance, and local NameCert rows.}
  \formalstatus{\unformalizedV}
  \bridgestatus{none}
  \notclaimed{This chapter does not close selected Lebesgue numbers, sequential compactness, host modulus functions, Lean verification, or bridge checking.}
  \upgradepath{Obligation closure requires checked source lock, finite-net coverage, pointwise-modulus coverage, rational lower-bound exactness, uniform-modulus output, and ledger non-escape for BEDC.Derived.CompactUniformModulusUp.}
  \constructivestory{A $\CompactUniformModulusUp$ packet is generated from a MetricUp source, CompactUp finite-net rows, ContinuousUp pointwise rows, a local finite fold ledger, RatUp comparisons, a UniformModulusUp output, componentwise hsame transport, Cont replay, provenance, and local naming.}
\end{closurestatus}
